\section{从历史缺陷到可数可加性}\label{sec:02-01}

\subsection{Riemann 与 Jordan 理论的边界}\label{subsec:2-1-1}
\index{Riemann 积分}\index{Jordan 内容}

Riemann 积分把一个区间切成有限多段，在每段上用一个矩形估计函数图像。这个办法的力量和
限制来自同一个词：有限。有限网格善于处理边界规整的对象，却没有任何机制把可数多次微小误差
合并起来。下面的例子分别提出“哪些函数可以积分”“哪些集合可以测量”和
“哪些极限可以穿过积分号”。

十九世纪末到二十世纪初出现了两条互补路线。一条从集合长度出发，先扩大可测集合，再把函数拆成水平集；
Lebesgue 理论沿这条路把可数可加性放在中心。另一条由 Young、Daniell 等人的工作推动，先研究连续函数
或简单函数上的正线性泛函，再从单调极限恢复更大的积分。下面暂不采用第二条路线的抽象定理，但两条路线
共同指出：有限分割本身不足以控制可数极限。

\begin{definition}[Darboux 上下和]\label{def:2-1-1-1}
设 $f:[a,b]\to\R$ 有界，分割
$P=\{a=x_0<x_1<\cdots <x_N=b\}$。记
\[
 M_j=\sup_{x\in[x_{j-1},x_j]}f(x),\qquad
 m_j=\inf_{x\in[x_{j-1},x_j]}f(x),
\]
并定义
\[
 U(f,P)=\sum_{j=1}^N M_j(x_j-x_{j-1}),\qquad
 L(f,P)=\sum_{j=1}^N m_j(x_j-x_{j-1}).
\]
上积分是所有 $U(f,P)$ 的下确界，下积分是所有 $L(f,P)$ 的上确界。二者相等时，称
$f$ Riemann 可积，其公共值记为 $\int_a^b f(x)\,\dd x$。
\end{definition}

\begin{example}[Dirichlet 函数]\label{ex:2-1-1-1}
Dirichlet 函数 $d=\1_{\Q\cap[0,1]}$ 在每个非退化区间内同时取到 $0$ 和 $1$。所以对每个分割
$P$，
\[
 U(d,P)=1,\qquad L(d,P)=0.
\]
它不 Riemann 可积。另一方面，$d$ 取值 $1$ 的点只有可数多个；若一种面积理论允许可数个点
仍然没有面积，那么它理应给 $d$ 积分 $0$。
\end{example}

\begin{example}[Thomae 函数]\label{ex:2-1-1-2}
在 $[0,1]$ 上定义
\[
 t(x)=
 \begin{cases}
 q^{-1},&x=p/q\text{，其中 }p\in\Z, q\in\N_{>0},\ \gcd(|p|,q)=1,\\
 0,&x\notin\Q.
 \end{cases}
\]
无理数在每个区间中稠密，故每个下和为 $0$。给定 $\varepsilon>0$，取 $N$ 使
$1/N<\varepsilon/2$。分母不超过 $N$ 的既约有理数在 $[0,1]$ 中只有有限多个，记为
$r_1,\dots,r_s$。对每个 $j$ 取 $r_j\in(\alpha_j,\beta_j)$，并使
\[
 \sum_{j=1}^s(\beta_j-\alpha_j)<\frac{\varepsilon}{2}.
\]
令 $P$ 包含 $0,1$ 以及所有落在 $[0,1]$ 内的 $\alpha_j,\beta_j$。把 $P$ 的小区间分成两类：
内部与 $\bigcup_j[\alpha_j,\beta_j]$ 相交于正长度的那些，和其余小区间。第一类小区间的长度和
不超过 $\sum_j(\beta_j-\alpha_j)$；第二类小区间的内部不含任何 $r_j$，故在其上
$t\le 1/N<\varepsilon/2$（分割端点不影响上确界，因为每个 $r_j$ 位于所选区间的内部）。因此
\[
 U(t,P)
 \le \sum_{j=1}^s(\beta_j-\alpha_j)
      +\frac1N\sum_{I\in P}|I|
 <\frac{\varepsilon}{2}+\frac{\varepsilon}{2}=\varepsilon.
\]
因此 $t$ Riemann 可积且积分为 $0$。它在每个有理点不连续、在每个无理点连续；与 Dirichlet
函数比较可见，坏点是否稠密并不是决定因素，坏点集合的“大小”才是。

最后两项连续性也可直接核对。若 $x\notin\Q$ 且 $\eta>0$，取 $N$ 使 $1/N<\eta$。分母不超过
$N$ 的既约有理数只有有限多个，且都不等于 $x$；取 $\delta>0$ 小于 $x$ 到这有限集的距离。于是
$|y-x|<\delta$ 时，要么 $y$ 无理而 $t(y)=0$，要么 $y$ 的既约分母大于 $N$ 而
$0<t(y)<\eta$，故 $t(y)\to t(x)=0$。若 $x=p/q$ 为既约有理数，取无理数列
$y_k\to x$，则 $t(y_k)=0$，而 $t(x)=1/q>0$，所以 $t$ 在 $x$ 不连续。
\end{example}

为了把集合本身纳入这项比较，我们先把有限盒计算说清楚。以下各盒一律采用半开约定
$\prod_{j=1}^n[a_j,b_j)$；改变有限个端面不影响本节的结论。

\begin{lemma}[初等集体积的良定义性]\label{lem:jordan-elementary-volume-well-defined}
若同一个有界集合 $A\subset\R^n$ 分别写成有限个两两不交半开盒的并
$A=\bigsqcup_{i=1}^r Q_i=\bigsqcup_{j=1}^s R_j$，则
\[
 \sum_{i=1}^r |Q_i|=\sum_{j=1}^s|R_j|,
 \qquad |\textstyle\prod_k[a_k,b_k)|:=\prod_k(b_k-a_k).
\]
\end{lemma}

\begin{proof}
收集所有 $Q_i,R_j$ 在第 $k$ 个坐标方向出现的有限多个端点，并按这些端点切分一条包含全部盒的
大区间。各坐标切分的笛卡尔积给出有限共同网格。每个 $Q_i$ 和 $R_j$ 都是若干网格小盒的
不交并。记共同网格胞腔为 $(G_\alpha)_\alpha$。一维恒等式
\[
 b-a=\sum_{l=1}^N(c_l-c_{l-1})
 \quad(a=c_0<\cdots<c_N=b)
\]
逐坐标相乘并展开，给出
\[
 |Q_i|=\sum_{G_\alpha\subset Q_i}|G_\alpha|,
 \qquad
 |R_j|=\sum_{G_\alpha\subset R_j}|G_\alpha|.
\]
两种不交表示所包含的网格胞腔都恰为满足 $G_\alpha\subset A$ 的那些，故
\[
 \sum_i|Q_i|
 =\sum_{G_\alpha\subset A}|G_\alpha|
 =\sum_j|R_j|.
\]
\end{proof}

\begin{definition}[Jordan 可测与内容]\label{def:2-1-1-2}
有限个两两不交半开盒的并称为\emph{初等集}，其体积由上一引理定义。对有界
$E\subset\R^n$，令
\[
 J^*(E)=\inf\{|A|:A\supset E,\ A\text{ 为初等集}\},\qquad
 J_*(E)=\sup\{|B|:B\subset E,\ B\text{ 为初等集}\}.
\]
若 $J^*(E)=J_*(E)$，称 $E$ Jordan 可测，公共值记作 $J(E)$。
\end{definition}

\begin{example}[可数并超出 Jordan 定义域]\label{ex:2-1-1-3}
单点 $\{x\}$ 可放在边长任意小的盒中，所以 $J^*(\{x\})=0$。然而
$\Q\cap[0,1]$ 是这些单点的可数并，并且它的闭包及其补集在 $[0,1]$ 中的闭包都是 $[0,1]$。
下面的边界判据将给出 $\Q\cap[0,1]$ 不 Jordan 可测。失败的不是单点计算，而是 Jordan 理论的
定义域不承受可数并。
\end{example}

\begin{example}[移动尖峰]\label{ex:2-1-1-4}
例如取归一化帐篷函数
\[
 \phi(u)=
 \begin{cases}
  16(u-\tfrac14),&\tfrac14\le u\le\tfrac12,\\
  16(\tfrac34-u),&\tfrac12<u\le\tfrac34,\\
  0,&\text{其余}.
 \end{cases}
\]
它非负、属于 $C_c((0,1))$，并且两个三角形各有底 $1/4$、高 $4$，故
$\int_0^1\phi(u)\,\dd u=2\cdot\frac12\cdot\frac14\cdot4=1$。在 $[0,1]$ 上令
$f_n(x)=n\phi(nx)$，并把 $\phi$ 在 $\R\setminus(0,1)$ 上延为 $0$。每个 $f_n$ 连续。若
$x>0$，当 $n>1/x$ 时 $nx>1$，故 $f_n(x)=0$；在 $x=0$ 也有 $f_n(0)=0$。所以
$f_n(x)\to0$ 处处成立，而换元 $u=nx$ 给出
\[
 \int_0^1 f_n(x)\,\dd x=\int_0^n\phi(u)\,\dd u=1.
\]
逐点收敛本身不足以交换极限和积分，而且失败即使在连续函数列中也会发生。
\end{example}

\begin{remark}[可数稳定性的几个侧面]\label{def:2-1-1-3}
上述例子把新理论所需的稳定性分成几个彼此独立的方面：非负单调极限与适当受控极限下的连续性，
不交可数分解上的可加性，平移和线性变换下的协变性，以及由连续函数或几何简单集给出的逼近。
这些性质将在不同定理中分别建立；它们并非一组等价条件。
\end{remark}

\begin{proposition}[Jordan 边界判据]\label{proposition:2-1-1-2}
有界集合 $E\subset\R^n$ Jordan 可测，当且仅当 $J^*(\partial E)=0$。此时，对任一包含
$E$ 的闭盒 $Q$，指标函数 $\1_E:Q\to\R$ Riemann 可积，且
\[
 J(E)=\int_Q\1_E(x)\,\dd x.
\]
\end{proposition}

\begin{proof}
先设 $J^*(\partial E)=0$。若 $\partial E=\varnothing$，则不存在同时遇到 $E$ 与 $E^c$ 的线段：
否则对这样一条线段的参数集合
\[
 T=\{t\in[0,1]:(1-t)x+ty\in E\}\qquad(x\in E,\ y\notin E)
\]
取 $T$ 与其补集的分界点，就会得到 $\partial E$ 中的点。因此任意网格上 $\1_E$ 都没有振荡胞腔，
以下结论立即成立。现在设 $\partial E\ne\varnothing$。给定 $\varepsilon>0$，由外内容为零，先取
有限半开盒 $Q_1,\dots,Q_r$ 覆盖 $\partial E$，使
\[
 \sum_{j=1}^r|Q_j|<\frac{\varepsilon}{2}.
\]
把每个 $Q_j$ 略微外扩成开盒 $U_j\supset Q_j$，并逐个控制
$|U_j|<|Q_j|+\varepsilon/(2r)$，于是
\[
 \partial E\subset O:=\bigcup_{j=1}^rU_j,
 \qquad \sum_{j=1}^r|U_j|<\varepsilon.
\]
$E$ 有界，故 $\partial E$ 紧；闭集 $O^c$ 与紧集 $\partial E$ 不交，所以
\[
 \delta:=\dist(\partial E,O^c)>0.
\]
把 $Q$ 切成直径小于 $\delta$ 的有限网格。若一个闭网格小盒 $C$ 同时遇到 $E$ 与
$Q\setminus E$，连接两点的线段包含某个 $z\in\partial E\cap C$。于是对每个 $w\in C$，
$|w-z|<\delta$，故 $w\notin O^c$，即 $C\subset O$。因此所有使 $\1_E$ 振幅为 $1$ 的小盒的并
包含于 $O$；用共同端点网格比较有限盒体积，得到这些小盒的总体积不超过
$\sum_j|U_j|$。于是
\[
 U(\1_E,P)-L(\1_E,P)\le\sum_j|U_j|<\varepsilon.
\]
于是 $\1_E$ Riemann 可积。更显式地，令 $\mathscr P$ 为该有限网格的胞腔族，并置
\[
 B_P=\bigcup_{C\in\mathscr P:\,\inf_C\1_E=1}C,
 \qquad
 A_P=\bigcup_{C\in\mathscr P:\,\sup_C\1_E=1}C.
\]
网格面的有限重叠为零内容，不影响这些初等集的体积，因而
\[
 B_P\subset E\subset A_P,qquad
 |B_P|=L(\1_E,P),\qquad |A_P|=U(\1_E,P).
\]
所以
\[
 0\le |A_P|-|B_P|=U(\1_E,P)-L(\1_E,P)<\varepsilon.
\]
这逐式给出 Jordan 内外逼近。由 Jordan 内容的定义以及 Darboux 可积性，分别有
\[
 |B_P|\le J(E)\le |A_P|,
 \qquad
 |B_P|=L(\1_E,P)\le\int_Q\1_E\le U(\1_E,P)=|A_P|.
\]
所以
\[
 \left|J(E)-\int_Q\1_E\right|
 \le |A_P|-|B_P|<\varepsilon;
\]
令 $\varepsilon\downarrow0$ 即得 $J(E)=\int_Q\1_E$。

反过来，设 $E$ Jordan 可测。取初等集 $B\subset E\subset A$ 且
$|A|-|B|<\varepsilon$，并取两者的共同网格。若 $x\in\partial E$ 落在某个网格小盒的内部，那么该
小盒若包含于 $B$，则 $x$ 是 $E$ 的内点；若不包含于 $A$，则 $x$ 是 $E^c$ 的内点。因此该小盒
必包含在 $A\setminus B$。这些胞腔两两不交，其总体积不超过 $|A|-|B|<\varepsilon$。网格边界是
有限个集合
\[
 \{x\in Q:x_i=c\}\qquad(1\le i\le n, c\text{ 为某个网格坐标}),
\]
每个都可用法向厚度为 $2\delta$ 的有限薄盒覆盖；有限求和后其总成本随 $\delta\downarrow0$ 趋于
零。若 $H$ 表示网格边界、$C_1,\dots,C_s$ 表示上述位于 $A\setminus B$ 的胞腔，则
\[
 \partial E\subset H\cup\bigcup_{l=1}^sC_l,
 \qquad
 J^*(\partial E)
 \le J^*(H)+\sum_{l=1}^s|C_l|
 \le J^*(H)+|A|-|B|.
\]
令 $\delta\downarrow0$ 得 $J^*(\partial E)\le\varepsilon$，再令 $\varepsilon\downarrow0$，得到
$J^*(\partial E)=0$。
\end{proof}

\begin{proposition}[Jordan 内容的有限性边界]\label{proposition:2-1-1-3}
Jordan 内容对有限个两两不交 Jordan 可测集可加，但 Jordan 可测集不构成 $\sigma$-代数。
具体地，每个单点内容为零，而 $\Q\cap[0,1]$ 的边界是 $[0,1]$，因而不 Jordan 可测。
\end{proposition}

\begin{proof}
若 $E_1,\dots,E_r$ 两两不交且 Jordan 可测，固定误差 $\varepsilon>0$，分别以初等集
$B_j\subset E_j\subset A_j$ 夹逼，使 $|A_j|-|B_j|<\varepsilon/r$。共同网格细分后，
$\bigsqcup_jB_j\subset\bigsqcup_jE_j\subset\bigcup_jA_j$；外并的体积至多各 $|A_j|$ 之和，
内并的体积等于各 $|B_j|$ 之和。于是
\[
 \sum_{j=1}^r|B_j|
 \le J_*\!\left(\bigsqcup_{j=1}^rE_j\right)
 \le J^*\!\left(\bigsqcup_{j=1}^rE_j\right)
 \le\left|\bigcup_{j=1}^rA_j\right|
 \le\sum_{j=1}^r|A_j|.
\]
两端之差小于 $\varepsilon$；令 $\varepsilon\downarrow0$ 得
$J(\bigsqcup_jE_j)=\sum_jJ(E_j)$。第二个断言由单点计算和上一命题得到。
\end{proof}

\begin{proposition}[单调极限迫使零集可数稳定]\label{proposition:2-1-1-4}
设 $I$ 定义在一个非负函数锥上，该锥含
$\1_{E_k},\1_{\cup_{k\le n}E_k},\1_{\cup_kE_k}$，对有限非负线性组合和有界递增极限封闭。
若 $I$ 可加、非负齐次、保持次序，并满足 $g_n\uparrow g$ 时 $I(g_n)\uparrow I(g)$，则
\[
 I(\1_{E_k})=0\quad(k\ge1)
 \quad\Longrightarrow\quad I(\1_{\cup_kE_k})=0.
\]
\end{proposition}

\begin{proof}
对每个 $n$，点态不等式
$\1_{\cup_{k\le n}E_k}\le\sum_{k=1}^n\1_{E_k}$ 与正性、有限线性给出
\[
 0\le I(\1_{\cup_{k\le n}E_k})
 \le\sum_{k=1}^nI(\1_{E_k})=0.
\]
左端指标函数随 $n$ 递增至 $\1_{\cup_kE_k}$，再用单调极限性质得到结论。
\end{proof}

\begin{theorem}[Lebesgue 的 Riemann 判据]\label{theorem:2-1-1-1}
有界函数 $f:[a,b]\to\R$ Riemann 可积，当且仅当它的不连续点集合具有 Lebesgue 测度零。
\end{theorem}

Lebesgue 测度将在\cref{subsec:2-1-3}中构造。判据的证明还需要函数振幅与微分理论，见
\cref{theorem:4-1-2-3}。

\begin{figure}[htbp]
\centering
\begin{tikzpicture}[node distance=13mm and 22mm,
  box/.style={draw,rounded corners,align=center,minimum width=31mm,minimum height=12mm},
  bad/.style={draw,dashed,rounded corners,align=center,inner sep=4pt},
  arr/.style={-{Latex[length=2mm]},thick}]
 \node[box] (j) {有限分割\\Jordan 内容};
 \node[box,right=of j] (l) {可数集合运算\\Lebesgue 测度};
 \node[box,right=of l] (d) {单调极限\\积分};
 \node[bad,below=of j] (q) {稠密可数并};
 \node[bad,below=of l] (dir) {Dirichlet 函数};
 \node[bad,below=of d] (sp) {移动尖峰};
 \draw[arr] (q)--(j);
 \draw[arr] (dir)--(l);
 \draw[arr] (sp)--(d);
 \draw[arr] (j)--node[midway,above,font=\small,fill=white,inner sep=1pt]{扩大定义域}(l);
 \draw[arr] (l)--node[midway,above,font=\small,align=center,fill=white,inner sep=1pt]{建立\\极限定理}(d);
\end{tikzpicture}
\caption{有限分割、可数集合运算与极限过程之间的三种失配。}
\label{fig:jordan-lebesgue-daniell-failures}
\end{figure}

\begin{remark}
Jordan 证明中真正失效的步骤是“取一个共同有限网格”。可数多个单点各自拥有小盒，不会自动产生
一个仍可由有限网格控制的覆盖。下一小节允许可数覆盖，并把近优覆盖的误差写成可求和的序列。
移动尖峰则只排除了“无控制的逐点收敛足以交换积分与极限”；一致收敛是一种修补，却不是唯一修补，第~3~章
还会建立单调收敛和受控收敛的不同条件。
\end{remark}

这些失败并不限于教材中的反例。Fourier 级数逼近和概率中的随机变量极限都天然先给逐点或几乎处处
信息；若积分要在这些极限下稳定，就必须补上可核查的控制条件。Young--Daniell 路线把这种压力写成
正泛函的单调连续性，集合路线则要求零集对可数并稳定。\cref{proposition:2-1-1-4} 已证明两者之间
最基本的联系；完整积分定理仍留到第~3~章。

\begin{exercise}
把上述 Thomae 函数证明改写成 Darboux 振幅判据的形式，并证明它恰在无理点连续。
\end{exercise}
\begin{exercise}
证明有限集合 Jordan 内容为零，并用边界判据再次证明 $\Q\cap[0,1]$ 不 Jordan 可测。
\end{exercise}
\begin{exercise}
构造一列连续非负函数，逐点趋于零但积分恒为给定常数 $c>0$。
\end{exercise}
\begin{exercise}
构造 Riemann 可积函数列，其逐点极限为 Dirichlet 函数。
\end{exercise}

有限网格已经无法同时追踪可数覆盖。下一小节保留“简单集合的覆盖成本”，但把覆盖数目放宽为可数，并
用一个切割条件从定义在所有子集上的外部价格中筛出真正可加的集合。

\subsection{外测度与 \Caratheodory{} 可测性}\label{subsec:2-1-2}
\index{外测度}\index{Caratheodory@\Caratheodory{} 可测性}

Jordan 理论试图从集合内部和外部同时有限夹逼。现在换一个次序：先不问集合是否可测，只允许用
简单集合从外面覆盖，并把所有覆盖成本的下确界当作外部价格。这个价格天然服从可数次可加，却
通常不能在所有集合上相加；可测性将由“切开任何测试集合都不损失成本”产生。

Lebesgue 的原始构造围绕实线上区间长度展开；\Caratheodory{} 的形式把其中真正使用的机制抽成外测度与
无损切割条件。这样，Euclidean 体积、Hausdorff 尺度和由预测度生成的测度可以共享一套证明。可测类
不是事先指定的全部子集，而是由“对每个测试集切割都不增加成本”挑出的、与该外测度相容的自然定义域。

在引入外部价格以前，先固定“测度”本身的含义。上一章已经定义了 $\sigma$-代数；现在增加的要求是
不交可数分解上的精确相加。

\begin{definition}[测度与测度空间]\label{def:2-1-2-measure-space}
设 $\mathcal A$ 是集合 $X$ 上的 $\sigma$-代数。函数
$\mu:\mathcal A\to[0,\infty]$ 称为\emph{测度}，若 $\mu(\varnothing)=0$，并且对任意两两不交的
$E_k\in\mathcal A$ 都有
\[
 \mu\Bigl(\bigsqcup_{k=1}^{\infty}E_k\Bigr)
 =\sum_{k=1}^{\infty}\mu(E_k).
\]
三元组 $(X,\mathcal A,\mu)$ 称为测度空间。若 $\mu(X)<\infty$，称 $\mu$ 为有限测度；若
$X=\bigcup_kX_k$ 且 $X_k\in\mathcal A$、$\mu(X_k)<\infty$，称 $\mu$ 为 $\sigma$-有限测度；
若 $\mu(X)=1$，称其为概率测度。拓扑空间上定义于 $\Borel(X)$ 的测度称为 Borel 测度。
性质 $P(x)$ 称为在 $\mu$-\emph{几乎处处}成立，若存在 $N\in\mathcal A$ 满足 $\mu(N)=0$，且
$P(x)$ 对每个 $x\in X\setminus N$ 成立；在未完成的测度空间中，这一定义不要求例外点集本身可测。
\end{definition}

可数可加性不只计算不交并；它还控制单调集合列的极限。后文的正则逼近和测度唯一性都会直接使用
下面两个公式。

\begin{proposition}[测度的单调连续性]\label{proposition:2-1-2-measure-continuity}
设 $(X,\mathcal A,\mu)$ 为测度空间。
\begin{enumerate}
 \item 若 $E_n\uparrow E$，则 $\mu(E_n)\uparrow\mu(E)$；
 \item 若 $E_n\downarrow E$ 且 $\mu(E_1)<\infty$，则 $\mu(E_n)\downarrow\mu(E)$。
\end{enumerate}
\end{proposition}

\begin{proof}
对递增列令 $D_1=E_1$、$D_n=E_n\setminus E_{n-1}$（$n\ge2$）。这些集合两两不交，且
$E=\bigsqcup_nD_n$、$E_N=\bigsqcup_{n\le N}D_n$，所以
\[
 \mu(E_N)=\sum_{n\le N}\mu(D_n)
 \uparrow\sum_{n=1}^{\infty}\mu(D_n)=\mu(E).
\]
若 $E_n\downarrow E$ 且 $\mu(E_1)<\infty$，则 $E_1\setminus E_n\uparrow E_1\setminus E$。
由第一项和有限可加性，
\[
 \mu(E_1)-\mu(E_n)=\mu(E_1\setminus E_n)
 \uparrow\mu(E_1\setminus E)=\mu(E_1)-\mu(E).
\]
这里所有被减数都不超过有限数 $\mu(E_1)$，故没有出现 $\infty-\infty$；移项即得第二项。
\end{proof}

\begin{example}[有限原子上的直接核算]\label{ex:2-1-2-finite-measure}
在 $X=\{1,\ldots,r\}$ 上取 $\mathcal A=2^X$，给定权重 $a_j\ge0$，并置
\[
 \mu(E)=\sum_{j\in E}a_j.
\]
任意不交可数族中至多有 $r$ 个非空成员，故可数可加性退化为有限和的重排，$\mu$ 是有限测度；当
$\sum_ja_j=1$ 时它是概率测度。这个模型说明测度公理怎样把有限权重相加，而外测度接下来要解决的
问题是：尚未给出可测 $\sigma$-代数时，怎样先为任意子集规定一致的外部价格。
\end{example}

\begin{definition}[外测度]\label{def:2-1-2-1}
集合 $X$ 上的函数 $\mu^*:2^X\to[0,\infty]$ 称为外测度，如果
\begin{enumerate}
 \item $\mu^*(\varnothing)=0$；
 \item $A\subset B$ 蕴含 $\mu^*(A)\le\mu^*(B)$；
 \item 对任意集合列 $(E_k)$，
 \(
 \mu^*(\bigcup_kE_k)\le\sum_k\mu^*(E_k).
 \)
\end{enumerate}
\end{definition}

\begin{definition}[覆盖生成的外测度]\label{def:2-1-2-3}
给定简单集合族 $\mathcal C\subset2^X$ 和成本 $c:\mathcal C\to[0,\infty]$，允许有限或可数覆盖，
并约定空族以成本 $0$ 覆盖空集，定义
\[
 \mu^*(E)=\inf\left\{\sum_{j=1}^{N}c(C_j):
 E\subset\bigcup_{j=1}^{N}C_j,\ C_j\in\mathcal C,\ 1\le N\le\infty\right\}.
\]
若不存在这样的覆盖，下确界取 $\infty$。
\end{definition}

\begin{lemma}[覆盖下确界给出外测度]\label{lemma:2-1-2-1}
若成本非负且空集有零成本空覆盖，则上一式定义外测度。
\end{lemma}

\begin{proof}
空集条件由约定得到。若 $A\subset B$，每个覆盖 $B$ 的族也覆盖 $A$，故下确界满足单调性。
只需证明可数次可加。若某个 $\mu^*(E_k)=\infty$，所需不等式的右端为无穷。否则，给定
$\varepsilon>0$，为每个 $k$ 选择一个有限或可数覆盖
$(C_{k,j})_{1\le j\le N_k}$（$N_k\le\infty$），使
\[
 \sum_{j=1}^{N_k}c(C_{k,j})<\mu^*(E_k)+\varepsilon2^{-k}.
\]
当 $E_k=\varnothing$ 时允许取空覆盖。所有这些有限或可数族的并仍是可数族，并覆盖
$\bigcup_kE_k$。于是
\[
 \mu^*\Bigl(\bigcup_kE_k\Bigr)
 \le\sum_{k,j}c(C_{k,j})
 \le\sum_k\mu^*(E_k)+\varepsilon.
\]
令 $\varepsilon\downarrow0$。这里同时为可数多个 $k$ 选择近优覆盖使用上一章
定义~\ref{def:1-4-3-1} 中的可数选择；证明并未假设任何下确界达到。
\end{proof}

\begin{example}[三种外部价格]\label{ex:2-1-2-1}
在 $\R^n$ 中以半开盒为 $\mathcal C$、以盒体积为成本，所得外测度记作 $m^*$。若
$E=\{x_1,x_2,\dots\}$ 可数，给第 $k$ 个点取体积小于 $\varepsilon2^{-k}$ 的盒，便有
$m^*(E)\le\varepsilon$，所以 $m^*(E)=0$。

另一个例子是计数外测度
\[
 \#^*(E)=\begin{cases}|E|,&E\text{ 有限},\\ \infty,&E\text{ 无限}.
 \end{cases}
\]
它稍后会成为所有集合均可测的情形。若改以直径的 $s$ 次幂为小尺度覆盖成本，则得到
Hausdorff 外测度；下一例说明这种成本如何反映集合的尺度。
\end{example}

\begin{example}[Hausdorff 外测度的覆盖动机]\label{ex:2-1-2-2}
若只用直径不超过 $\delta$ 的集合覆盖 $E$，并把成本取作
$\sum_j(\diam U_j)^s$，把一条长度为 \(r\) 的线段等分为 \(N\) 段便得到总成本
\[
 N\left(\frac rN\right)^s=r^sN^{1-s}.
\]
在 \(s=1\) 时该成本恒为 \(r\)，在 \(s>1\) 时则随 \(N\to\infty\) 趋于零；单点在每个
\(s>0\) 时也有零成本。先固定 $\delta$ 取覆盖下确界，再令 $\delta\downarrow0$，便得到能分辨尺度的
外部价格。\Cref{subsec:2-4-3} 将给出正式定义及 Lipschitz 映射下的覆盖畸变不等式。
\end{example}

\begin{definition}[\Caratheodory{} 可测]\label{def:2-1-2-2}
给定外测度 $\mu^*$，集合 $E\subset X$ 称为 $\mu^*$-可测，如果对每个 $A\subset X$，
\[
 \mu^*(A)=\mu^*(A\cap E)+\mu^*(A\setminus E).
 \label{eq:caratheodory-cut}
\]
\end{definition}

外测度次可加已经给出等式中“$\le$”的一边；定义要求切开后的总成本不超过整体最优成本。

\begin{example}[可测性依赖外测度]\label{ex:2-1-2-3}
对计数外测度，若 $A$ 有限，等式 \eqref{eq:caratheodory-cut} 是有限集合的基数分解；若 $A$ 无限，
右侧至少一项为无穷，等式仍成立。因此每个集合都可测。
\end{example}

\begin{example}[不是每个外测度都测量所有集合]\label{ex:2-1-2-4}
设 $X$ 至少有两个点，令 $\mu^*(\varnothing)=0$，每个非空集合的外测度为 $1$。它满足外测度
三条公理。若 $E$ 为真非空子集，以 $A=X$ 测试，则
\[
 1=\mu^*(X)\ne\mu^*(E)+\mu^*(X\setminus E)=2.
\]
所以 $E$ 不可测。这说明 \Caratheodory{} 可测性不是集合脱离外测度后的绝对属性。
\end{example}

\begin{proposition}[只需有限成本测试]\label{proposition:2-1-2-4}
检验 $E$ 的 \Caratheodory{} 条件时，只需考虑 $\mu^*(A)<\infty$ 的测试集。
\end{proposition}

\begin{proof}
当 $\mu^*(A)=\infty$ 时，次可加性给
$\infty\le\mu^*(A\cap E)+\mu^*(A\setminus E)$，所以右端也是无穷，等式自动成立。
\end{proof}

\begin{theorem}[\Caratheodory{} 定理]\label{theorem:2-1-2-2}
所有 $\mu^*$-可测集构成一个 $\sigma$-代数 $\mathcal M$，并且
$\mu=\mu^*|_{\mathcal M}$ 是完备测度。特别地，若 $N\in\mathcal M$ 且 $\mu(N)=0$，则
$N$ 的每个子集都属于 $\mathcal M$ 且测度为零。
\end{theorem}

\begin{proof}
空集满足切割等式，补集稳定来自等式对 $E$ 和 $E^c$ 的对称性。设 $E,F$ 可测。对任意测试集
$A$，先以 $E$ 切割，再以 $F$ 切割 $A\setminus E$，得到
\[
\begin{aligned}
 \mu^*(A)
 &=\mu^*(A\cap E)+\mu^*(A\setminus E)\\
 &=\mu^*(A\cap E)+\mu^*((A\setminus E)\cap F)
   +\mu^*(A\setminus(E\cup F))\\
 &\ge \mu^*(A\cap(E\cup F))+\mu^*(A\setminus(E\cup F)).
\end{aligned}
\]
最后一步使用
$A\cap(E\cup F)=(A\cap E)\cup((A\setminus E)\cap F)$ 的次可加性。反向不等式由把 $A$ 分成
这两块的次可加性得到，所以 $E\cup F$ 可测；归纳得到有限并稳定。

现在设 $E_k\in\mathcal M$，令
$D_1=E_1$、$D_k=E_k\setminus\bigcup_{j<k}E_j$。有限并和补集稳定保证 $D_k$ 可测且两两
不交。记 $D=\bigcup_kD_k$。对任意测试集 $A$，逐次用 $D_1,\dots,D_n$ 切割，得到
\[
 \mu^*(A)\ge
 \sum_{k=1}^n\mu^*(A\cap D_k)
 +\mu^*\Bigl(A\setminus\bigcup_{k=1}^nD_k\Bigr).
 \label{eq:caratheodory-iterate}
\]
最后一项不小于 $\mu^*(A\setminus D)$。令 $n\to\infty$，
\[
 \mu^*(A)\ge\sum_{k=1}^{\infty}\mu^*(A\cap D_k)+\mu^*(A\setminus D)
 \ge\mu^*(A\cap D)+\mu^*(A\setminus D),
\]
第二个不等式是外测度次可加。反向不等式同样来自次可加，故 $D$ 可测。这证明 $\mathcal M$
为 $\sigma$-代数。

在 \eqref{eq:caratheodory-iterate} 中取 $A=D$ 并令 $n\to\infty$，与次可加的反向估计合并，
得到 $\mu(D)=\sum_k\mu(D_k)$，所以限制是测度。

最后，设 $N$ 可测且 $\mu^*(N)=0$，并取 $S\subset N$。对任意 $A$，单调性给
$\mu^*(A\cap S)=0$。又因为
$A=(A\setminus S)\cup(A\cap S)$，
\[
 \mu^*(A)\le\mu^*(A\setminus S)+0\le\mu^*(A).
\]
所以 $S$ 满足 \Caratheodory{} 等式，并且零测。
\end{proof}

\begin{definition}[度量外测度]\label{def:2-1-2-4}
度量空间上的外测度 $\mu^*$ 若在 $\dist(A,B)>0$ 时满足
$\mu^*(A\cup B)=\mu^*(A)+\mu^*(B)$，称为\emph{度量外测度}（metric outer measure）。
\end{definition}

\begin{proposition}[度量外测度的 Borel 性]\label{proposition:2-1-2-3}
度量外测度的每个闭集都 \Caratheodory{} 可测，因而每个 Borel 集可测。
\end{proposition}

\begin{proof}
设 $F$ 闭。由 \cref{proposition:2-1-2-4}，只需取 $\mu^*(A)<\infty$。令
\[
 C_j=A\cap\{2^{-j}\le\dist(x,F)<2^{-j+1}\}\quad(j\ge1),
\]
并把 $\dist(x,F)\ge1$ 的部分记为 $C_0$。若 $k\ge j+2$，则对
$x\in C_j,y\in C_k$ 有
\[
 \dist(x,F)-\dist(y,F)
 \ge2^{-j}-2^{-k+1}\ge2^{-j-1}.
\]
距离函数是 $1$-Lipschitz，故 $d(x,y)\ge2^{-j-1}>0$。因此同为奇指标或同为偶指标的任意有限
子族彼此可逐项使用度量可加性，给出
\[
 \sum_{j\le N,\ j\text{ odd}}\mu^*(C_j)\le\mu^*(A),\qquad
 \sum_{j\le N,\ j\text{ even}}\mu^*(C_j)\le\mu^*(A).
\]
令 $N\to\infty$ 后相加，得到
\[
 \sum_{j=1}^{\infty}\mu^*(C_j)\le2\mu^*(A)<\infty,
\]
所以该级数收敛。

令 $B_n=A\cap\{\dist(x,F)\ge2^{-n}\}$。由
\[
 A\setminus F\subset B_n\cup\bigcup_{j>n}C_j
\]
和外测度次可加性，
\[
 \mu^*(A\setminus F)
 \le \mu^*(B_n)+\sum_{j>n}\mu^*(C_j).
 \label{eq:metric-outer-ring-tail}
\]
另一方面，$B_n$ 与 $A\cap F$ 正距离，故
\[
 \mu^*(A)\ge\mu^*(B_n\cup(A\cap F))
 =\mu^*(B_n)+\mu^*(A\cap F).
\]
集合 $B_n$ 递增，故数列 $b_n=\mu^*(B_n)$ 递增；又 $B_n\subset A$ 且 $\mu^*(A)<\infty$，
所以 $b_n\uparrow b<\infty$。在 \eqref{eq:metric-outer-ring-tail} 中令 $n\to\infty$，环带级数的
尾和趋零，得到 $\mu^*(A\setminus F)\le b$。另一方面，上式的正距离加性对每个 $n$ 给
$b_n+\mu^*(A\cap F)\le\mu^*(A)$；令 $n\to\infty$ 后合并两式，得到
\[
 \mu^*(A)\ge\mu^*(A\setminus F)+\mu^*(A\cap F).
\]
反向由次可加，故 $F$ 可测。闭集生成 Borel $\sigma$-代数，结论随之成立。
\end{proof}

\begin{definition}[测度空间的完成化]\label{def:measure-completion}
设 $(X,\mathcal A,\mu)$ 是测度空间，令
\[
 \begin{aligned}
 \overline{\mathcal A}^{\mu}
 &=\{E\subset X:E\mathbin\triangle A\subset Z
       \text{ 对某些 }A,Z\in\mathcal A,\ \mu(Z)=0\}\\
 &=\{A\cup N:A\in\mathcal A,\ N\subset Z\in\mathcal A,\ \mu(Z)=0\}.
 \end{aligned}
\]
若 $E\mathbin\triangle A$ 包含于可测零集，置 $\bar\mu(E)=\mu(A)$，称所得测度空间为原空间的
完成化。
\end{definition}

\begin{proposition}[完成化的代表]\label{prop:measure-completion-representative}
$\overline{\mathcal A}^{\mu}$ 是 $\sigma$-代数，上述 $\bar\mu$ 良定义且完备。每个完成化可测集
$E$ 都有 $A\in\mathcal A$ 使 $E\mathbin\triangle A$ 包含于一个 $\mu$-零集。
\end{proposition}

\begin{proof}
若 $E=A\cup N$，则 $E\mathbin\triangle A=N\setminus A$ 包含于可测零集。反过来，若
$E\mathbin\triangle A\subset Z$ 且 $Z$ 可测零，则
$E=(A\setminus Z)\cup(E\cap Z)$，所以定义中的两种表示等价。补集和可数并时，可测代表分别取补集和
可数并，误差仍包含于可数个可测零集之并，故该类为 $\sigma$-代数。若两个表示给出代表 $A,B$，
则 $A\mathbin\triangle B$ 包含于可测零集，于是
$\mu(A\setminus B)=\mu(B\setminus A)=0$，从可加性得到 $\mu(A)=\mu(B)$；故 $\bar\mu$ 良定义。

设两两不交的 $E_i$ 有可测代表 $A_i$，且 $E_i\mathbin\triangle A_i\subset Z_i$，其中
$\mu(Z_i)=0$。同时选取可数族 $(A_i,Z_i)$ 是一次可数选择；在通常的 \textsf{ZFC} 背景下这是
允许的。令
\[
 B_1=A_1,\qquad B_i=A_i\setminus\bigcup_{j<i}A_j\quad(i\ge2).
\]
则 $B_i$ 可测且两两不交。由于 $E_i\cap E_j=\varnothing$，有
$A_i\cap A_j\subset Z_i\cup Z_j$；从而 $A_i\mathbin\triangle B_i$ 包含于有限个零集之并，
故 $B_i$ 仍是 $E_i$ 的代表。因此
\[
 \left(\bigcup_iE_i\right)\mathbin\triangle
 \left(\bigcup_iB_i\right)
 \subset\bigcup_i(E_i\mathbin\triangle B_i)
\]
包含于可测零集 $\bigcup_iZ_i'$，其中每个 $Z_i'$ 是相应的有限零集并；故
$\bigsqcup_iB_i$ 确实是 $\bigsqcup_iE_i$ 的可测代表。于是
\[
 \bar\mu\Bigl(\bigsqcup_iE_i\Bigr)
 =\mu\Bigl(\bigsqcup_iB_i\Bigr)
 =\sum_i\mu(B_i)=\sum_i\bar\mu(E_i).
\]
这证明可数可加性。若 $\bar\mu(E)=0$，取零测可测代表 $A$；则 $E$ 包含于 $A$ 与表示误差零集之并，
所以 $E$ 的任意子集仍属于完成化且测度为零。最后一句已包含在上述对称差表示中。
\end{proof}

\begin{proposition}[完成化可测函数的代表]\label{prop:completion-measurable-function-representative}
设 $\bar\mu$ 是 $(X,\mathcal A,\mu)$ 的完成化。若
$f:X\to[-\infty,\infty]$ 满足
$f^{-1}(B)\in\overline{\mathcal A}^{\mu}$ 对每个 Borel 集
$B\subset[-\infty,\infty]$ 成立，则存在具有
$g^{-1}(B)\in\mathcal A$ 的 $g:X\to[-\infty,\infty]$，使 $f=g$ 在某个 $\mu$-零集之外成立。
\end{proposition}

\begin{proof}
先设 $0\le f\le1$。令
\[
 s_n=2^{-n}\lfloor 2^nf\rfloor,
\]
并在 $f=1$ 处令 $s_n=1$。$s_n$ 只取有限多个值；它的每个水平集属于完成化，因而有
$\mathcal A$ 中的代表。具体地，把 $s_n$ 的不同值列为 $v_{n,0},\dots,v_{n,r_n}$，令
$L_{n,j}=\{s_n=v_{n,j}\}$，并取 $A_{n,j}\in\mathcal A$ 与可测零集 $Z_{n,j}$，使
\[
 L_{n,j}\mathbin\triangle A_{n,j}\subset Z_{n,j}.
\]
指标对 $(n,j)$ 的集合可数；同时作这些选择使用可数选择。按固定次序不交化：
\[
 B_{n,0}=A_{n,0},\qquad
 B_{n,j}=A_{n,j}\setminus\bigcup_{l<j}A_{n,l}\quad(1\le j\le r_n),
\]
并在余集 $X\setminus\bigcup_jB_{n,j}$ 上指定默认值 $v_{n,0}$。令 $t_n=v_{n,j}$ 于
$B_{n,j}$。这是 $\mathcal A$-可测有限值函数。若
$x\notin Z_n:=\bigcup_jZ_{n,j}$，则 $x$ 对所有 $A_{n,j}$ 的隶属关系与对 $L_{n,j}$ 的隶属关系
相同；水平集 $(L_{n,j})_j$ 恰好分割 $X$，故此时恰有一个代表含 $x$，且
$t_n(x)=s_n(x)$。

置 $Z=\bigcup_nZ_n$ 和 $g=\limsup_nt_n$。为显式核对可测性，对每个 $a\in\R$ 有
\[
 \{g>a\}
 =\bigcup_{\substack{q\in\Q\\q>a}}
   \bigcap_{N=1}^{\infty}\bigcup_{n\ge N}\{t_n>q\}\in\mathcal A.
\]
这些半直线生成扩展实线的 Borel $\sigma$-代数，故 $g$ 是 $\mathcal A$-可测的。在
$X\setminus Z$ 上，
$t_n=s_n\to f$，所以 $g=f$。

一般情形取把扩展实线双射到 $[0,1]$ 的严格递增连续函数，例如在有限点置
$\psi(x)=\tfrac12+\pi^{-1}\arctan x$，并令 $\psi(-\infty)=0,\psi(\infty)=1$。对
$\psi\circ f$ 使用上一步，再与 Borel 可测的 $\psi^{-1}$ 复合。
\end{proof}

\begin{figure}[htbp]
\centering
\begin{tikzpicture}[scale=.95]
 \path[fill=BookAccent!10]
   (-3,-1.6)--(0,-1.6)..controls(-1,-.5)and(1,.5)..(0,1.6)--(-3,1.6)--cycle;
 \path[fill=BookWarm!12]
   (0,-1.6)--(3,-1.6)--(3,1.6)--(0,1.6)..controls(1,.5)and(-1,-.5)..(0,-1.6)--cycle;
 \draw[rounded corners,thick] (-3,-1.6) rectangle (3,1.6);
 \draw[BookInk,thick] (0,-1.6)..controls(-1,-.5)and(1,.5)..(0,1.6);
 \node at (-1.65,0) {$A\cap E$};
 \node at (1.65,0) {$A\setminus E$};
 \node[above] at (0,1.6) {测试集 $A$};
 \draw[-{Latex[length=2mm]},thick] (0,-2.15)--(0,-1.55);
 \node[below,align=center] at (0,-2.15)
 {$\mu^*(A)=\mu^*(A\cap E)+\mu^*(A\setminus E)$\\表示分别覆盖不增加最优成本};
\end{tikzpicture}
\caption{\Caratheodory{} 条件把任意测试集无损切成两块。}
\label{fig:caratheodory-lossless-cut}
\end{figure}

\begin{remark}
外测度与测度不可混同。外测度定义在所有子集上，却通常不加；真正的测度是它在可测 $\sigma$-代数
上的限制。相反，完成化是在已经有测度之后，把可测零集的所有子集补入定义域。
\end{remark}

覆盖外测度将在本章产生 Lebesgue 测度，在\cref{subsec:2-4-3}产生 Hausdorff 尺度，也可用于路径
空间或乘积空间上的外部构造。另一条存在性路线是先在环上给预测度，再作 \Caratheodory{} 扩张；它在
\cref{subsec:2-2-1}正式建立。两条路线都用可数覆盖，但“从所有子集中筛选”与“从简单集合向外延拓”
承担不同任务。

\begin{exercise}
从覆盖定义逐项验证 Lebesgue 外测度的平移不变性。
\end{exercise}
\begin{exercise}
证明零外测度集合满足 \Caratheodory{} 条件。
\end{exercise}
\begin{exercise}
在 $|X|\ge2$ 时给出一个外测度，使恰有 $\varnothing$ 与 $X$ 可测。
\end{exercise}
\begin{exercise}
补写度量外测度命题中环带内缩和奇偶尾估计的全部细节，并据此证明开集可测。
\end{exercise}

外测度给出了抽象筛选机制；要把它变成 Euclidean 体积，还需从盒覆盖中证明平移、缩放以及开紧逼近，
而不能把这些几何性质当作定义的一部分。

\subsection{Lebesgue 测度：Euclidean 空间中的正则体积}\label{subsec:2-1-3}
\index{Lebesgue 测度}\index{正则性}

现在把覆盖成本取成 Euclidean 盒体积。仅有可数可加性还不足以刻画体积：它还应尊重
平移与缩放，并允许用开集从外、紧集从内逼近。这些性质来自盒几何、紧性和
Euclidean 空间的可数拓扑骨架。

Lebesgue 理论的关键进展，是先建立一个对可数集合运算稳定的 Borel 测度，
再把 Borel 零集的所有子集补入完成化。完成化扩大了可测对象，却保留几何体积；Jordan 内容和 Riemann
积分则应当作为规整边界上的旧理论被重新识别，而不是与新理论平行共存、互不相干。

\begin{definition}[Lebesgue 外测度与测度]\label{def:2-1-3-1}
对半开盒 $Q=\prod_{j=1}^n[a_j,b_j)$ 置 $|Q|=\prod_j(b_j-a_j)$，并令
\[
 m^*(E)=\inf\left\{\sum_k|Q_k|:E\subset\bigcup_kQ_k\right\}.
\]
把 $m^*$ 的 \Caratheodory{} 可测类暂记为 $\mathcal L_C$，限制测度暂记为 $m_C$。
\end{definition}

\begin{definition}[Borel 与 Lebesgue $\sigma$-代数]\label{def:2-1-3-2}
$\Borel(\R^n)$ 是开集生成的 $\sigma$-代数。下文将先证明
$\Borel(\R^n)\subset\mathcal L_C$；把 $m_C$ 限制到 $\Borel(\R^n)$ 得 Borel Lebesgue 测度
$m_B$。$m_B$ 的完成化记为 $\mathcal L$，称为 Lebesgue $\sigma$-代数，其测度记作 $m$。
\Cref{lem:lebesgue-completion-characterization} 将证明 $\mathcal L_C=\mathcal L$ 且
$m_C=m$；在那以前两种构造的记号保持分开。
\end{definition}

\begin{definition}[正则性]\label{def:2-1-3-3}
拓扑空间上的测度 $\mu$ 称外正则，如果
\[
 \mu(E)=\inf\{\mu(U):E\subset U,\ U\text{ 开}\};
\]
称在 $E$ 上内正则，如果
\[
 \mu(E)=\sup\{\mu(K):K\subset E,\ K\text{ 紧}\}.
\]
\end{definition}

\begin{lemma}[盒的外测度]\label{lem:lebesgue-box-lower-bound}
每个半开盒 $Q$ 满足 $m^*(Q)=|Q|$。
\end{lemma}

\begin{proof}
单盒覆盖给 $m^*(Q)\le|Q|$。若某条边长为零，则 $Q=\varnothing$（按半开盒约定），等式成立。
以下设各边长 $d_i=b_i-a_i>0$。取
\[
 K_\eta=\prod_{i=1}^n[a_i+\eta,b_i-\eta]
 \subset Q,
 \qquad0<\eta<\tfrac12\min_i d_i.
\]
给定可数半开盒
覆盖 $(Q_j)$；若 $\sum_j|Q_j|=\infty$，所需下界平凡。以下设总成本有限。给定 $\varepsilon>0$，
把 $Q_j$ 向外略微放大为开盒 $U_j$，使 $|U_j|<|Q_j|+\varepsilon2^{-j}$。紧集 $K_\eta$ 有有限子覆盖
$U_{j_1},\dots,U_{j_r}$。把 $K_\eta$ 与这些有限开盒在各坐标方向的所有端点组成共同网格。每个正体积
网格胞腔 $C$ 若包含于 $K_\eta$，取 $z\in C^\circ$。某个 $U_{j_l}$ 包含 $z$；因为网格包含
$U_{j_l}$ 的全部端点，所以 $C^\circ\subset U_{j_l}$。把每个胞腔分配给一个这样的 $l$。
胞腔边界位于有限个坐标超平面中；用法向总厚度小于 $\tau$ 的有限薄盒覆盖这些超平面后，有限盒
体积的可加性给
\[
 |K_\eta|\le\sum_{l=1}^r|U_{j_l}|+\tau
 \le\sum_j|Q_j|+\varepsilon+\tau.
\]
先令 $\tau\downarrow0$、再令
$\varepsilon\downarrow0$，得到
\[
 |K_\eta|=\prod_{i=1}^n(d_i-2\eta)\le\sum_j|Q_j|.
\]
最后令 $\eta\downarrow0$，得到
$|Q|\le\sum_j|Q_j|$。对所有覆盖
取下确界即得结论。
\end{proof}

\begin{lemma}[任意集合的开外逼近]\label{lem:lebesgue-outer-open-approximation}
对任意 $A\subset\R^n$ 与 $\varepsilon>0$，存在开集 $U\supset A$ 使
\[
 m^*(U)\le m^*(A)+\varepsilon.
\]
当 $m^*(A)<\infty$ 时还可取严格不等号。
\end{lemma}

\begin{proof}
若 $m^*(A)=\infty$，取 $U=\R^n$ 即可。以下设 $m^*(A)<\infty$。取半开盒覆盖
$A\subset\bigcup_iQ_i$，使
\[
 \sum_i|Q_i|<m^*(A)+\frac{\varepsilon}{2}.
\]
对每个 $i$ 选开盒 $O_i$ 与半开盒 $R_i$，使
\[
 Q_i\subset O_i\subset R_i,
 \qquad |R_i|<|Q_i|+\varepsilon2^{-i-1}.
\]
令 $U=\bigcup_iO_i$。则 $U$ 开、$A\subset U\subset\bigcup_iR_i$，从而
\[
 m^*(U)\le\sum_i|R_i|
 <\sum_i|Q_i|+\frac{\varepsilon}{2}
 <m^*(A)+\varepsilon.
\]
\end{proof}

\begin{proposition}[外测度可由任意细立方体计算]\label{proposition:2-1-3-1a}
给定 $\delta>0$，$m^*(E)$ 等于所有边长不超过 $\delta$ 的可数半开立方体覆盖的总体积下确界。
\end{proposition}

\begin{proof}
立方体是盒，所以允许所有盒所得下确界不大于所述下确界。若 $m^*(E)=\infty$，反向不等式平凡。
以下设 $m^*(E)<\infty$，并取总成本小于 $m^*(E)+\varepsilon/2$ 的盒覆盖 $(R_k)$。对每个 $k$，
选 $0<h_k\le\delta$ 足够小，并用网格 $h_k\Z^n$ 中所有与 $R_k$ 相交的半开立方体覆盖 $R_k$。
若 $R_k$ 的边长为 $d_{k,1},\ldots,d_{k,n}$，这些立方体的总体积至多
\[
 \prod_{j=1}^n(d_{k,j}+2h_k),
\]
它在 $h_k\downarrow0$ 时趋于 $|R_k|$。故可令该总量小于
$|R_k|+\varepsilon2^{-k-1}$。合并所有覆盖得到总成本小于 $m^*(E)+\varepsilon$，再令
$\varepsilon\downarrow0$。
\end{proof}

\begin{remark}
薄盒 $[0,1]\times[0,\varepsilon]$ 面积趋零而直径趋近 $1$。因此以后用 Lipschitz 映射控制像的
直径时，不能直接传送任意盒覆盖；必须先使用上一命题把覆盖立方体化。
\end{remark}

\begin{lemma}[Lebesgue 外测度的度量可加性]\label{lem:lebesgue-metric-outer}
若 $\dist(A,B)>0$，则
\[
 m^*(A\cup B)=m^*(A)+m^*(B).
\]
因此 $m^*$ 是度量外测度，所有 Borel 集都属于 $\mathcal L_C$；特别地，每个半开盒
\Caratheodory{} 可测。
\end{lemma}

\begin{proof}
次可加性给“$\le$”。设 $\delta=\dist(A,B)>0$，取边长小于 $\delta/\sqrt n$ 的固定网格。
给定 $A\cup B$ 的可数半开盒覆盖，把每个覆盖盒沿该网格切成至多可数个不交半开小盒；小盒总体积
等于原盒体积，每个小盒的直径小于 $\delta$，故不可能同时遇到 $A$ 与 $B$。收集所有遇到 $A$ 的
小盒得到 $A$ 的覆盖，收集所有遇到 $B$ 的小盒得到 $B$ 的覆盖；两族的总体积之和不超过原覆盖的
总成本。因此
\[
 m^*(A)+m^*(B)\le\sum_k|Q_k|
\]
对 $A\cup B$ 的每个盒覆盖 $(Q_k)$ 成立。取下确界得到“$\ge$”。于是 $m^*$ 是 metric outer
measure，\cref{proposition:2-1-2-3} 给出 Borel 可测性。半开盒是 Borel 集，最后一句随之成立。
\end{proof}

\begin{theorem}[Lebesgue 测度的基本几何性质]\label{theorem:2-1-3-1}
Lebesgue 测度在 $\R^n$ 上平移不变；对 $r>0$ 和 Lebesgue 可测 $E$，
\[
 m(rE)=r^n m(E).
\]
每个半开盒 $Q$ 满足 $m(Q)=|Q|$。
\end{theorem}

\begin{proof}
平移一个盒不改变成本，所以逐项平移覆盖给 $m^*(E+x)=m^*(E)$；用 $-x$ 得反向不等式。
对缩放，任一盒覆盖 $(Q_k)$ 被送到盒覆盖 $(rQ_k)$，其成本为
$\sum_k|rQ_k|=r^n\sum_k|Q_k|$；对逆缩放再做一次，故
$m^*(rE)=r^nm^*(E)$。这些等式表明平移和缩放保持 \Caratheodory{} 等式。它们还是
Euclidean 同胚，因而保持 Borel 集；外测度公式又说明它们把 Borel 零集送到 Borel 零集，所以也
保持完成化 $\mathcal L$。盒公式由\cref{lem:lebesgue-box-lower-bound,lem:lebesgue-metric-outer}
得到。
\end{proof}

\begin{theorem}[正则性与 $\sigma$-有限性]\label{theorem:2-1-3-2}
$\R^n=\bigcup_{k\ge1}[-k,k]^n$，每个耗尽盒有限测。Borel 集外正则，开集内正则；Lebesgue
完成化对每个可测集外正则，并对每个有限测度可测集内正则。紧集具有有限测度。
\end{theorem}

\begin{proof}
若 $U$ 开，对每个 $x\in U$ 可取有理端点闭盒 $K\subset U$ 且 $x\in\interior K$。
这样的盒只有可数多个，故 $U=\bigcup_jK_j$。令 $L_N=\bigcup_{j\le N}K_j$；它紧且
$L_N\uparrow U$。测度的可数可加性蕴含从下连续，故
$m(U)=\sup_Nm(L_N)$，得到开集内正则。

设 $E$ 为 Borel 集。若 $m(E)=\infty$，则每个包含 $E$ 的开集也有无穷测度，外正则等式平凡。
以下设 $m(E)<\infty$。由引理 \ref{lem:lebesgue-outer-open-approximation}，对每个 $j$ 可取
开集 $U_j\supset E$，使
\[
 m(U_j)=m^*(U_j)<m^*(E)+2^{-j}=m(E)+2^{-j}.
\]
令 $j\to\infty$，得到 Borel 集外正则。

若 $E$ 是有限测度 Borel 集，由 $[-R,R]^n\uparrow\R^n$ 的从下连续性，先取 $R$ 使
$m(E\setminus[-R,R]^n)<\varepsilon/2$。对 Borel 集 $[-R,R]^n\setminus E$ 作外逼近，取开集
$V$ 包含它并满足
\[
 m\bigl(V\setminus([-R,R]^n\setminus E)\bigr)<\varepsilon/2.
\]
令 $F=[-R,R]^n\setminus V$。则 $F$ 是包含于 $E$ 的紧集，而且
\[
 E\setminus F
 \subset (E\setminus[-R,R]^n)
 \cup\bigl(V\setminus([-R,R]^n\setminus E)\bigr),
\]
故 $m(E\setminus F)<\varepsilon$。这给有限 Borel 集内正则。

完成化可测集 $E$ 有 Borel 代表 $B$ 且 $E\mathbin\triangle B$ 零测。若 $m(E)=\infty$，外正则
等式平凡。若 $m(E)<\infty$，取 Borel 零集 $Z\supset E\mathbin\triangle B$。给定
$\varepsilon>0$，由上一引理取开集 $O\supset Z$ 使 $m(O)<\varepsilon/2$，并由 Borel 外正则性
取开集 $V\supset B$ 使 $m(V)<m(B)+\varepsilon/2$。则 $E\subset V\cup O$ 且
\[
 m(V\cup O)\le m(V)+m(O)<m(E)+\varepsilon,
\]
故 $E$ 外正则。继续使用上述 Borel 零集 $Z$。先取紧集 $K\subset B$，使
$m(B\setminus K)<\varepsilon/2$；
再取开集 $O\supset Z$，使 $m(O)<\varepsilon/2$。于是 $K\setminus O$ 是包含于 $E$ 的紧集，并且
\[
 m\bigl(E\setminus(K\setminus O)\bigr)
 \le m(E\setminus B)+m(B\setminus K)+m(K\cap O)<\varepsilon.
\]
这证明完成化的有限测度集合内正则。对于无限测度集合，内逼近须先在有限盒中局部化；不能把两个无穷
测度作差来制造误差估计。
\end{proof}

\begin{proposition}[Borel 代表]\label{proposition:2-1-3-3}
每个 Lebesgue 可测集 $E$ 存在 $G_\delta$ 集 $G\supset E$ 和 $F_\sigma$ 集 $F\subset E$，使
$m(G\setminus E)=m(E\setminus F)=0$。
\end{proposition}

\begin{proof}
令 $Q_k=[-k,k]^n$，并置 $D_1=Q_1$、$D_k=Q_k\setminus Q_{k-1}$（$k\ge2$）。这些 Borel 层
两两不交、覆盖 $\R^n$，且每层有限测。对每个 $j,k\ge1$，由外正则性取开集
$U_{j,k}\supset E\cap D_k$，使
\[
 m\bigl(U_{j,k}\setminus(E\cap D_k)\bigr)<2^{-j-k}.
\]
令 $U_j=\bigcup_kU_{j,k}$、$G=\bigcap_jU_j$。则 $G$ 是 $G_\delta$ 集且包含 $E$。又因为
\[
 U_j\setminus E
 \subset\bigcup_k\bigl(U_{j,k}\setminus(E\cap D_k)\bigr),
\]
所以 $m(U_j\setminus E)\le2^{-j}$，进而 $m(G\setminus E)=0$。

再由有限测度集合的内正则性，对每个 $k$ 取紧集 $K_k\subset E\cap Q_k$，使
$m((E\cap Q_k)\setminus K_k)<2^{-k}$，并令 $F=\bigcup_kK_k$。这是包含于 $E$ 的
$F_\sigma$ 集。固定 $k$ 后，对每个 $\ell\ge k$ 都有
\[
 (E\cap Q_k)\setminus F
 \subset (E\cap Q_\ell)\setminus K_\ell,
\]
故其测度不超过 $2^{-\ell}$；令 $\ell\to\infty$，得到
$m((E\cap Q_k)\setminus F)=0$。最后对 $k$ 取可数并，便有 $m(E\setminus F)=0$。
\end{proof}

\begin{proposition}[两种 Lebesgue 可测类相同]\label{lem:lebesgue-completion-characterization}
$\mathcal L_C=\mathcal L$，并且 $m_C=m$。
\end{proposition}

\begin{proof}
完成化包含于 $\mathcal L_C$，因为该类含 Borel 集且 $m_C$ 完备。反过来，设
$E\in\mathcal L_C$，把空间分为不交有限测度盒层 $D_k$。对每个 $j,k$，先取开集
$U_{j,k}\supset E\cap D_k$；由引理 \ref{lem:lebesgue-outer-open-approximation} 可使
$m^*(U_{j,k})<m^*(E\cap D_k)+2^{-j-k}$。集合 $E\cap D_k$ 是 \Caratheodory{} 可测的且
外测度有限；以它切割可测开集 $U_{j,k}$ 得
\[
 m^*(U_{j,k})=m^*(E\cap D_k)+m^*(U_{j,k}\setminus(E\cap D_k)),
\]
所以差集外测度小于 $2^{-j-k}$。令
$U_j=\bigcup_kU_{j,k}$ 和 $G=\bigcap_jU_j$。则 $G$ 为 Borel 集，$E\subset G$，并且
$m^*(G\setminus E)\le2^{-j}$ 对每个 $j$ 成立，故差集外测度为零。令 $N=G\setminus E$。
对每个 $r\ge1$，由 $m^*(N)=0$ 取开集 $O_r\supset N$ 且 $m(O_r)<2^{-r}$，再置
\[
 Z=\bigcap_{r=1}^{\infty}O_r.
\]
则 $Z$ 是包含 $N$ 的 Borel 集，并且对每个 $r$ 都有
$0\le m(Z)\le m(O_r)<2^{-r}$，故 $m(Z)=0$。因此 $E$ 属于
Borel 完成化，且两种限制测度都等于 $m^*(E)$。
\end{proof}

\begin{example}[可数集零测度]\label{ex:2-1-3-1}
前面的逐点覆盖证明给出每个可数集零测。特别地，Dirichlet 函数取值 $1$ 的集合虽在每个区间
稠密，却在 Lebesgue 意义下可忽略。
\end{example}

\begin{example}[Cantor 集]\label{ex:2-1-3-2}
三分 Cantor 集 $C$ 在第 $k$ 阶包含于 $2^k$ 个长度 $3^{-k}$ 的闭区间并中，故
$m(C)\le2^k3^{-k}=(2/3)^k$。令 $k\to\infty$ 得 $m(C)=0$。它却不可数、紧且无内点，说明基数、
拓扑大小和测度大小彼此不同。
\end{example}

\begin{example}[完成化确实增加集合]\label{ex:2-1-3-3}
Euclidean 空间的 Borel 集由一个可数基生成。把“取补集”和“取可数并”按可数序数逐层编码，每个
Borel 集可由一个可数码描述；所有可数码的基数至多为
$|\N^\N|=|\R|$，故 Borel 集总数至多为连续统。另一方面，映射
\[
 (\varepsilon_j)_{j\ge1}\longmapsto
 \sum_{j=1}^{\infty}2\varepsilon_j3^{-j},
 \qquad \varepsilon_j\in\{0,1\},
\]
把 $\{0,1\}^{\N}$ 双射到 Cantor 集 $C$（端点采用只含 $0,2$ 的三进制展开），所以
$|C|=|\R|$。Cantor 定理给 $|2^C|>|\R|$，故 $C$ 至少有一个非 Borel 子集 $A$。由于
$A\subset C$ 且
$m(C)=0$，完成化使 $A$ Lebesgue 可测且零测。
\end{example}

\begingroup
\raggedbottom
\begin{proposition}[Jordan 理论嵌入]\label{proposition:2-1-3-4}
有界集合 $E$ Jordan 可测，当且仅当 $m(\partial E)=0$；此时 $E$ Lebesgue 可测且
$J(E)=m(E)$。
\end{proposition}

\begin{proof}
若 $E$ Jordan 可测，\cref{proposition:2-1-1-2} 给出边界的有限任意小盒覆盖，故
$m(\partial E)=0$。反过来，$\partial E$ 是有界闭集，因而紧。给定 $\varepsilon>0$，先以总体积
小于 $\varepsilon/2$ 的可数半开盒 $(Q_j)$ 覆盖它，再把第 $j$ 个盒略微外扩为开盒 $O_j$，使
\[
 |O_j|<|Q_j|+\varepsilon2^{-j-2}.
\]
紧性选出有限子覆盖 $O_{j_1},\dots,O_{j_r}$。对每个选中开盒再取包含它的半开盒 $R_l$，并把
有限个增量分配成
\[
 \sum_{l=1}^r(|R_l|-|O_{j_l}|)<\frac{\varepsilon}{4}.
\]
于是
\[
 \sum_{l=1}^r|R_l|
 <\sum_{j=1}^{\infty}|Q_j|+\frac{\varepsilon}{4}+\frac{\varepsilon}{4}
 <\varepsilon.
\]
故 $J^*(\partial E)=0$，Jordan 边界判据给出
Jordan 可测性。

在任一方向，只要 $m(\partial E)=0$，就有
$E=\interior{E}\cup(E\setminus\interior{E})$，其中第一项为 Borel 集，第二项包含于零测闭集
$\partial E$；Lebesgue 测度的完备性因此给出 $E$ 可测。最后，给定 $\varepsilon>0$，取初等集
$B\subset E\subset A$ 使 $|A|-|B|<\varepsilon$。初等集的 Lebesgue 测度等于其体积，故
\[
 |B|=m(B)\le m(E)\le m(A)=|A|,
 \qquad
 |B|\le J(E)\le |A|.
\]
于是 $|m(E)-J(E)|\le |A|-|B|<\varepsilon$；令 $\varepsilon\downarrow0$ 即得
$J(E)=m(E)$。
\end{proof}

\begin{figure}[htbp]
\centering
\begin{tikzpicture}[scale=1]
 \draw[BookAccent,thick,rounded corners] (-3,-1.5) rectangle (3,1.5) node[above left] {$U_j$};
 \draw[BookInk,thick] plot[smooth cycle] coordinates {(-2,-.7)(-1.4,1)(-.2,.55)(.9,1.05)(2.1,.15)(1.5,-1)(.2,-.7)(-1.1,-1.1)};
 \node at (0,.05) {$E$};
 \draw[BookWarm!90,thick,rounded corners] (-1.2,-.55) rectangle (1.05,.55) node[below right] {$K_j$};
 \draw[<->,thin] (2.2,-1.45)--(2.2,-1.0) node[midway,right] {$\to0$};
 \draw[<->,thin] (-1.2,-.55)--(-1.2,-.95) node[midway,left] {$\to0$};
\end{tikzpicture}
\caption{Lebesgue 正则性：开集从外、紧集从内逼近。箭头标记的是测度误差，而非 Euclidean 距离。}
\label{fig:lebesgue-inner-outer-regularity}
\end{figure}

\begin{remark}
平移不变和单位盒正规化不能在 $2^{\R^n}$ 上产生可数可加体积；\Caratheodory{} 构造必须舍弃某些
子集。另一个常见混淆来自完成化：连续映射自动拉回 Borel 集，却不自动拉回任意完成化集合。
若要拉回完成化，需另证目标零测 Borel 集的逆像在源中零测，或者在拉回后重新完成。
\end{remark}

这一体积将成为第~3~章积分、换元与卷积的底层测度，并在 Fourier 分析中提供平移不变的积分背景。
正则性还会把集合逼近转化为 $C_c(\R^n)$ 的函数逼近；概率密度则是在这个测度之上描述分布的一种
方式。相应的积分、卷积与密度理论从下一章开始建立。

\begin{exercise}
证明中间三分 Cantor 集不可数、无内点并且零测。
\end{exercise}
\begin{exercise}
证明每个开集是可数个两两仅在边界相交的半开 dyadic 盒之并。
\end{exercise}
\begin{exercise}
从正则性重新证明每个 Lebesgue 可测集有 Borel 代表。
\end{exercise}
\begin{exercise}
先对任意 $A\subset\R^n$ 从覆盖定义证明
$m^*(A+x)=m^*(A)$ 与 $m^*(rA)=r^nm^*(A)$；再证明这些映射保持可测性，并对可测 $A$ 写成相应的
$m$ 等式。
\end{exercise}

下一节不再依赖 Euclidean 盒。我们将把“存在一个扩张”和“扩张由简单测试唯一确定”分开处理，
并把拓扑正则性、有符号抵消纳入同一语言。
\clearpage
\endgroup
